\section{Introduction}

The voltage--conductance equation describes a population of neurons
by a joint probability density in voltage and synaptic conductance.
Conductance diffusion supplies the noise, voltage transport carries
neurons to threshold, and a boundary reset returns a firing neuron
to its resting voltage. The firing flux also changes the conductance
drift and diffusion. The resulting autonomous equation can therefore
couple fluctuations of the population to its own firing rhythm.

Let $v_E>v_F>0$ and $y_L,y_*,a_*,c>0$. On
$D=(0,v_F)\times(0,\infty)$ put
\begin{align}\label{eq:full-coefficients}
 J(v,y)&=y(v_E-v)-y_Lv,&
 y_F&=\frac{y_Lv_F}{v_E-v_F},\notag\\
 K_F(t,y)&=y_*+cN(t)-y,& a_F(t)&=a_*+c^2N(t).
\end{align}
The firing rate is
$N(t)=\int_{y_F}^\infty J(v_F,y)F(t,v_F,y)\dd y$.
We seek a nonnegative mass-one density satisfying
\begin{equation}\label{eq:full-equation}
 \partial_tF+\partial_v(JF)+\partial_y(K_FF)-a_F\partial_{yy}F=0 .
\end{equation}
Writing $J_0(y)=J(0,y)$ and $J_F(y)=J(v_F,y)$, its voltage
traces obey
\begin{equation}\label{eq:full-voltage-boundaries}
 \begin{cases}
  F(t,0,y)=F(t,v_F,y)=0,&0<y<y_F,\\
  J_0(y)F(t,0,y)=J_F(y)F(t,v_F,y),&y>y_F.
 \end{cases}
\end{equation}
The conductance flux $Q=K_FF-a_FF_y$ has zero trace at $y=0$
and vanishes at infinity.
Voltage conditions are kinetic weak traces; the floor condition is
the natural conormal condition in the distributional formulation.
At infinity we mean that $Q$ paired against every smooth periodic
time test and smooth closed-voltage test has tail limit zero.
There is no prescribed value at the single grazing conductance $y_F$.

For this full equation, Fonte Sanchez and Mischler developed the
prescribed linear evolution and stationary theory, including
exponential relaxation in the frozen linear problem
\cite[Theorems 1.2 and 2.2]{FM24}. Fonte Sanchez, Mischler and Salort
proved global weak existence for general connectivity
\cite[Theorem 1.1 and Definition 4.1]{FMS26}. These results give
the evolution and trace framework used here. A nonconstant
periodic trajectory is a separate long-time question:
the voltage--conductance discussion in the survey of Carrillo
and Roux explicitly lists periodic-solution existence as open
\cite[Section 7.6]{CR25}. Carrillo, Dou and Zhou studied a
separable simplified equation and its fast-conductance limit;
their periodic solutions concern that limit
\cite[Section 4]{CDZ22}. The result below concerns the full
nonseparable transport field, positive conductance diffusion,
and both firing-feedback terms in \eqref{eq:full-coefficients}.

\begin{theorem}\label{thm:main}
There are positive constants $y_*,a_*,c,T$ such that, with
$v_E=2$, $v_F=1$ and $y_L=1$, equations
\eqref{eq:full-coefficients}--\eqref{eq:full-voltage-boundaries}
admit a nonnegative $T$-periodic distributional density $F$
of mass one whose firing rate $N$ is nonconstant.
The density is continuous in time into $L^1(D)$, has the
stated voltage and zero conductance-flux traces, and satisfies
\begin{equation}\label{eq:main-moment-entropy}
 \sup_{0\le t\le T}\int_D
       \bigl(y^2F(t,v,y)+|F(t,v,y)\log F(t,v,y)|\bigr)\dd v\dd y<\infty .
\end{equation}
One may choose
$F\in H^1_{\per}(0,T;L^2(D,e^{4y}\dd v\dd y))$ and
$N\in H^1_{\per}(0,T)$, with nonzero first Fourier coefficient
of $N$. The convention is $0\log0=0$.
\end{theorem}

The construction uses small frozen conductance diffusion and a mean
strictly above the deterministic firing threshold. One firing cycle
then provides a useful phase coordinate. We compare its noisy return
operator with an OU evolution over the deterministic cycle in a
common Gaussian weighted boundary space. The comparison holds in
operator norm for arbitrary data in that space. A corrected phase
mode, transferred through its isolated one-dimensional projection,
gives an actual simple characteristic root. Its real part has a
negative noise contribution and a positive mean-feedback contribution.
The reflecting-floor part of diffusion feedback is retained throughout.

Two further points are needed to obtain the theorem. First, the
fixed-noise physical resolvent is analytic in the mean and diffusion,
so a countable zero-set argument selects parameters with no other
integer harmonic resonance. The mean and coupling then serve as
two free parameters in a fixed-period rate equation. Second, its
initially signed density response is identified with the fixed
point of a positive prescribed period map. This step establishes
nonnegativity without requiring a pointwise positive lower bound
for the stationary density. Stability of the resulting periodic
solutions and explicit numerical parameter values are not determined
by this construction.

Section~\ref{sec:linear} constructs the divergence inverse with its
complete outgoing trace and natural floor input.
Section~\ref{sec:returns} establishes the process and cycle estimates.
Section~\ref{sec:crossing} proves the simple full-response crossing,
and Section~\ref{sec:periodic} closes the periodic rate equation and
proves the theorem. Appendix~\ref{app:boundary} gives the boundary
and process constructions; Appendix~\ref{app:clock} supplies the
bounded-clock comparison.

All complex transpose pairings below are bilinear; only energy
inner products use conjugation. We keep separate the random first
hit $\tf$, deterministic cycle time $\tau$, and final period $T$.
Constants in fixed-positive-noise resolvent estimates may depend on
that noise. Small-noise estimates are uniform only on the specified
compact mean and frequency sets, and harmonic expansions are for
each fixed harmonic.
